\section{投资问题：什么会改变中性／观察}

本报告的结论不是静态目标价，而是一组可证伪条件。核心判断是，恒逸风险分为盈利、现金和每股三类：价差与负荷决定利润，库存与资本开支决定现金，少数股东与股本决定每股归属。只要其中一类显著恶化，即使另外两类暂时良好，估值也可能下调。

“高风险”在本报告中指：单一变量的合理不利变动可使基准价值下降超过15\%，或需要外部融资才能维持既定扩张。文莱裂解、营运资金、高杠杆和资本开支属于高风险；可转债和少数股东属于中高每股风险；安全环保事件概率较低但影响高。

\begin{exhibitbox}[风险矩阵]
\scriptsize
\noindent\begin{tabularx}{\textwidth}{L{2.75cm}C{1.2cm}C{1.2cm}L{3.15cm}X}
\toprule
\textbf{风险} & \textbf{概率} & \textbf{影响} & \textbf{领先指标} & \textbf{触发动作} \\
\midrule
文莱裂解价差均值回归 & 高 & 高 & 柴油／航煤、PX价差、负荷 & 下调H2利润与PE \\
国内PTA／聚酯恢复不及预期 & 中高 & 中高 & PTA加工费、长丝库存、织机开工 & 下调国内分部与2027路径 \\
营运资金吞噬现金 & 中高 & 高 & OCF、库存、预付、短债 & PE向6倍收敛 \\
高杠杆与利息 & 高 & 高 & 净债务代理、利息覆盖、再融资 & 扩大资本折价 \\
文莱二期／煤化工超支 & 中 & 高 & 预算、融资、进度、担保 & 二期维持零信用，压低FCFE \\
少数股东分流 & 高 & 中 & 总净利与归母差额 & 下调归母转换率 \\
可转债与库存股回流 & 高 & 中 & 转股、赎回、股份用途 & 重算股本、现金与利息 \\
股份支付费用 & 中 & 中 & 过户、公允价值、摊销 & 下调报告口径利润 \\
汇率／原料／运费错配 & 中 & 中高 & USD／CNY、Brent、航运 & 调整单位毛利情景 \\
安全环保与停产 & 低中 & 高 & 事故、处罚、检修公告 & 事件性降级 \\
关联交易／担保 & 中 & 高 & 预付、应收、担保余额与代偿 & 降低质量评分与倍数 \\
\bottomrule
\end{tabularx}
\end{exhibitbox}
\sourcenote{analysis/risk\_framework.md；HY-OFF-2025AR；HY-OFF-2026Q1；研究截止2026-07-22。概率和影响为AStock House判断。}

风险之间会联动：裂解回落降低利润，库存和短债恶化现金，转股或员工计划扩大分母；三者同时发生时，目标价下行不是线性。尾部10.5元反映的是利润65亿元与6倍PE共同收缩，而非单一盈利偏差。

\section{催化剂：只有能改变分子、分母或倍数的事件才重要}

正面催化包括H1实际归母与扣非落在预告中上沿、经营现金流显著转正、文莱分部披露高负荷和可解释少数股东、国内PTA与聚酯库存下降、CPL--PA6达到盈亏平衡，以及二期融资条款透明。单纯股价上涨、行业宣传或新规划不构成基本面催化。

\begin{exhibitbox}[催化与证伪日历]
\small
\noindent\begin{tabularx}{\textwidth}{L{2.65cm}L{3.15cm}L{3.15cm}X}
\toprule
\textbf{窗口} & \textbf{正面催化} & \textbf{负面证伪} & \textbf{模型动作} \\
\midrule
2026H1半年报 & 归母55--60、扣非54.7--59.7亿元；OCF显著转正 & 低于预告或现金仍大额流出 & 更新全年利润与PE \\
2026Q3／三季报 & Q3单季约12亿元、9M累计约69.5亿元；文莱仍高负荷 & Q3单季低于8亿元或9M低于约65.5亿元、库存短债上升 & 重估Q4与全年情景 \\
7月28日后转债窗口 & 转股降债与利息同步下降 & 股份供给增加但降债不明显 & 更新分母、债务和利息 \\
员工计划过户 & 认购、现金、费用透明 & 费用高于缓冲或认购结构不利 & 更新利润和股本 \\
文莱二期融资／工程 & 金额、利率、资本金、提款与进度明确 & 融资或工期不确定性上升 & 决定是否给予期权信用 \\
国内链月度 & PTA加工费、长丝库存和织机开工共振 & 加工费回落、再累库 & 更新国内利润与营运资金 \\
\bottomrule
\end{tabularx}
\end{exhibitbox}
\sourcenote{HY-OFF-2026H1P；HY-OFF-CB；HY-OFF-ESOP7；HY-OFF-P2；analysis/chain\_earnings\_bridge.md，截止2026-07-22。}

催化剂必须与模型字段相连。现金流改善主要提高倍数，H2利润直接调整分子，转债和员工计划调整分母及现金，二期金融交割只提高远期期权。这样可以避免同一事件同时在利润、倍数和期权中重复计价。

\section{监控仪表盘：下一次更新看什么}

中报应完成四项核验：预告准确性、文莱与国内分部、少数股东归属、现金流与库存。三季报应完成H2利润坡度核验。月度行业数据只能提前预警，不能替代公司季报。可转债、员工计划和二期公告属于事件监控，应在发生时即时重算。

\begin{exhibitbox}[核心监控阈值]
\noindent\begin{tabularx}{\textwidth}{L{3.1cm}L{2.75cm}L{2.75cm}X}
\toprule
\textbf{指标} & \textbf{升级阈值} & \textbf{降级阈值} & \textbf{估值影响} \\
\midrule
H1实际归母／扣非 & 在预告区间且接近中上沿 & 低于55／54.7亿元 & 调整全年利润 \\
H1经营现金流 & 明显转正、库存下降 & 仍大额为负且库存继续升 & PE上调／下调 \\
Q3／9M归母 & 单季约12／累计约69.5亿元或更高 & 单季低于8／累计低于约65.5亿元 & 重估Q4；全年H2低于18亿元则基准失效 \\
少数股东占比 & 约30\%且可解释 & 显著更高且无解释 & 调整归母转换 \\
PTA加工费 & 高于2025均值且不累库 & 约250元／吨以下或再累库 & 国内利润调整 \\
聚酯库存／织机 & 库存下降、开工回升 & 库存上升、开工低于57\% & 毛利与营运资金调整 \\
净债务／短债 & 随转股和现金改善下降 & 存货、短债、在建同步升 & 资本折价调整 \\
二期融资与进度 & 完整条款、进度快于3.17\% & 融资／工期不透明 & 期权信用调整 \\
\bottomrule
\end{tabularx}
\end{exhibitbox}
\sourcenote{analysis/risk\_framework.md；analysis/chain\_earnings\_bridge.md；分析阈值截止2026-07-22。行业阈值为代理，不等于公司利润。}

\section{行动框架与最终观点}

13.5元以下且利润、裂解和现金没有证伪时，进入条件研究升级区；13.5--16.3元维持观察；16.3--16.5元降低研究优先级；16.5元以上若没有盈利和现金流上修则触发重新评估。12元以下若伴随基本面恶化，不作机械均值回归。该框架是研究纪律，不执行订单。

\begin{houseviewbox}[最终投资结论]
恒逸石化2026H1利润反转真实，文莱一期是核心驱动，国内链条处于不均衡修复；但当前15.06元已经计入已披露稀释口径下约97\%的House基准利润。模型值15.6575元，发布目标15.7元，核心公允区间13.5--16.5元，尾部10.5--22.5元，评级中性／观察。\textbf{下一次升级不看同比标题，看H2利润、经营现金流、少数股东和资本开支四项是否同时改善。}
\end{houseviewbox}

\begin{riskbox}[最终失效条件]
若H2归母低于18亿元、H1现金流不修复、文莱裂解显著回落或资本开支／融资压力上升，本报告下调；若H2达到35亿元以上且现金覆盖利息与资本开支、二期融资受控，本报告上调牛情景。剩余2.31亿库存股用途和员工计划股份支付费用在披露前始终是正式边界。
\end{riskbox}
